\chapter{Consolidated Discussion, Limits of Inference, and Conclusion}
\label{ch:discussion}\label{ch:conclusion}\label{chap:discussion}

The programme now supports an assessment across the completed mathematical, synthetic and regional empirical arc. Its central observed-catalogue result is the failure of the tested inhibition additions to improve matched predictive scores. Large spatial gains mainly correct coarse routing and uniform cell-lift representations; these are useful diagnostic checks, but do not establish competitive forecasting superiority. Standard unweighted inverse-power ETAS beats OC in every reported California and Japan primary fold. The tested inhibition additions remain losing or tied in the observed comparisons, including after temporal refinement and localization. These statements concern declared procedures and common scored observations, not a universal ranking of all inhibitory and excitation models.
\section{What the theory separates}

The mathematical results delimit several distinct regimes. Finite positive-mode systems admit the stated first-moment stability conclusions under their link, mark and subcritical-triggering assumptions. Finite-horizon approximation does not establish an interchange of stationary and infinite-mode limits. Exact-memory counterexamples apply to specified no-feedback submodels, and the linear-response calculation concerns an auxiliary frozen-slope system. Full-law identifiability also requires assumptions stronger than those available in an observed finite catalogue. The preceding mathematical chapters state their assumptions and proofs. The new comparison in Section~\ref{sec:autonomous-nz-comparison} directly fits an autonomous finite-mode subfamily satisfying the first-moment hypotheses. Its early and late folds select positive inhibition and lose held-out score; the middle fold selects $\eta=0$ and is tied at displayed precision. This connects the stability model to an executed fitted comparison while keeping stability and predictive advantage distinct.

This separation matters for the empirical discussion. A physical model may contain loading or initial-stress memory while a finite catalogue lacks power to distinguish it from a restricted alternative. A diagnostic family may improve its training score because its likelihood geometry favours extra structure, without thereby identifying the full physical law. A small held-out loss may be numerically controlled yet limited as statistical evidence when folds overlap or were inspected. Conversely, a non-rejection is not evidence of equality. Failure to distinguish two recipes leaves a set of explanations: the data may be insufficient, the candidate may be misspecified, the optimizer may be imperfect, or the tested effect may be small over the observed horizon. The controlled studies weaken some specific explanations, but do not select a unique physical one.

At \(\eta=0\), a separately varied stress kernel disappears from the diagnostic intensity when background and triggering are held fixed. It is therefore unidentified there. The analogous physical loading argument requires finite loading and initial-stress parameters: it cannot be carried into a singular small-feedback, large-loading limit without additional analysis. Non-rejection of a shared-kernel restriction at or near zero feedback consequently does not validate equality of two physical memories.

Synthetic diagnostics further clarify, but do not erase, the inferential gap. The historical comparison with nuisance quantities supplied produced 150 detections in 150 trials, whereas a fitted setting produced 57 in 150. These are different simple-versus-composite inferential tasks; their contrast is not an intrinsic information ceiling, and it is not proof that the remainder is avoidable algorithmic inefficiency. It offers no justification for fixing the real background merely to recover apparent power. Better optimization or a larger synthetic bank can improve numerical precision without changing the scientific question.

The transformed marks add another qualification. The magnitude excess has an exponential law, while the stress and productivity transforms have finite first and infinite second moments at the declared generating parameters. That does not prevent bounded-indicator Monte Carlo summaries, but it invalidates unqualified use of Fisher-information shortcuts based on second moments. Computational stability and statistical efficiency therefore have to be defended under the declared scoring rule, not inferred from formal differentiability alone.


\section{What the temporal and mark extensions resolve}

The magnitude-history and productivity studies retain negative matched results. A class-history law can change conditional mark probabilities and, through joint fitting, timing parameters without providing additional held-out information. Its effective-scale algebra concerns the fixed observed marked history; it does not establish equality of unconditional multistep future laws. Free productivity likewise changes the fitting problem and its finite-moment regime. Its failed continuation gate is part of the result, not an omitted successful analysis waiting to be inferred.

Exposure orthogonalization separates numerical geometry from the represented family. Better Gram conditioning is useful, but an interior optimum that maps feasibly into both coordinate boxes need not give a better predictive model. The negligible score change in Stage A is therefore consistent with the algebra. Its coarse Omori bridge must be distinguished from Extension O's continuously tuned temporal law. The latter improves excitation yet leaves matched and jointly tuned inhibition losing. The national temporal loss cannot be attributed solely to having fixed the original ML triggering kernel.

Common free ML time scale improves its own temporal comparison. Parent-class time scales add an inconsistent increment and fail the registered continuation rule, leaving observed association fitting unexecuted. This does not erase the earlier completed class-history experiment: the two branches have distinct references, parameters and gates. Every claim of nonexecution in this monograph refers to its specified branch.

\section{What localization changes, and what it does not}

Regional nonclosure gives a mathematical reason to test local structure: equal pooled depletion can conceal different sums of local rates. The completed progression in Chapter~\ref{ch:nz-spatial} supplies the empirical comparison that the earlier manuscript proposed. It moves from zone-local excitation to redistribution, distance routing, projected continuous density and a factorial with stronger immigrant geography. These changes improve particular location scores, but incremental inhibition remains negative, boundary-selected or tied. A reduced whole-model penalty after localization must not be confused with a positive inhibition-minus-matched-excitation contrast.

Continuous geography retains its advantage over the cell-lift representation under the richer background. Location dominates that benefit. Temporal and background contrasts still contain losses, including changes whose pooled joint sign differs from their timing sign and whose fold signs differ from their pooled sign. Spatial normalization leaves regional intensity unchanged at fixed shared parameters; timing changes after refitting do not imply that a normalized spatial density directly creates regional events. Nor can gains from separate historical stages be concatenated as if they were one fitted comparison.

The selected spatial scales describe a bounded statistical recipe in a declared projection and domain. Endpoint selection, negative small-scale sensitivities, repeated coordinates and unknown coordinate accuracy prevent physical interaction-distance claims. Training density is an estimate of event geography under the observed record, not an identified distribution of immigrants. Missing parents, region boundaries and catalogue selection remain observation-law limitations.

\section{Transfer, competition and conditional adequacy are different questions}

California retains the primary matched temporal and location directions, and the fitted target beats its no-trigger reference. EP nevertheless wins the joint comparison in every original primary window, chiefly through conditional location. The higher-magnitude sensitivity and annual/sequence accounting retain their separate denominators and do not convert retrospective inspection into untouched confirmation. The Ridgecrest active partition is a predetermined-in-form retrospective time union; it is not a causal assignment of individual aftershocks. These regional score conclusions use recovered original fit and score records with checked hashes and source identities, rather than newly recomputed regional scores; the original audit scopes remain distinct from this recovery.

Japan transfers the continuous-versus-lifted location direction but reverses the Omori-versus-fixed-ML timing direction in every scored window. The combined descriptive transfer rule therefore fails. EP again wins jointly. A statement that ``the improvement transfers'' without naming the component and reference would obscure a central negative finding. The explicit complete-primary selection rule also limits the population to which the JMA comparison applies; excluded unresolved records are not silently imputed.

The frozen OC conditional simulations ask another question. Material count, gap, spatial and subblock discrepancies survive the complete identity/count/mark/exposure audit. Tied diagnostic values require insertion intervals, particularly for empty cells; a single randomized rank would hide that ambiguity. The full raw replay has narrower path coverage than the aggregate audit. These diagnostics show material conditional inadequacy of OC across the three origins. Deterministic immigrant exposure is an accounting coordinate, not a random calibration test, and total/trigger exposure ranks are redundant. EP's log-score victory stands on its own; Section~\ref{sec:s05-completed} shows fixed-state discrepancies for EP too \citep{Schorlemmer2007,Zechar2010}.

The completed matched magnitude-history study (Section~\ref{sec:s02-completed}) separates frequent small positive fitted gains from evidence calibrated against its empirical null bank. Matching productivity reduces the mean GLR relative to the misspecified candidate, but positive selection remains common. Exceedances of the same-bank empirical threshold do not establish independent Type I error control or a physical explanation.

The four-arm spatial comparison (Section~\ref{sec:s03-completed}) addresses a different contrast from the earlier coarse cell-lift checks. EP retains a regional joint advantage over OC in California and New Zealand, chiefly through location. Fold contrasts vary; the comparison neither establishes universal EP superiority nor causally decomposes tail shape and background effects.

The expanded California backtest (Section~\ref{sec:s06-completed}) has a distinct retrospective episode roster and variable procedure contrasts: EP does not win every episode. Its causal training prefixes do not make the inspected roster prospective or independent. In Japan, fixing the background trend at zero reduces the pooled OC--FC timing deficit but leaves every fold negative. The trend ablation therefore does not remove the reversal or rescue the combined descriptive transfer rule.

\section{Continuous local inhibition and the zero-strength boundary}

The LI extension attenuates only the immigrant contribution without spatial renormalization. The field uses unit stress marks and all eligible strictly earlier recorded parents. Every eligible event therefore contributes the same inhibition amplitude regardless of magnitude, while triggering is unattenuated. This is a restrictive diagnostic of immigrant attenuation, not a physical stress-shadow model or whole-intensity suppression. At zero strength it nests OC exactly, while range and memory disappear from the law.

All selected observed California LI states lie at that boundary. The tiny retained positive floating-point difference from OC is a numerical tie, not a predictive gain. The separately archived OC--EP and LI--EP values retain distinct source identities even though their displayed rounded values agree. Range and memory display conventions at zero do not identify physical scales. EP's better score remains a competitive result rather than a calibration certificate.

Every fixed-seed regime matters. In the historical initial eight-dataset study, all four null controls lose held-out score, including a loss of 148.65 nats. This shows that nesting the null in training does not protect future performance: estimation variability and overfitting remain plausible. Held-out losses alone do not establish an optimization defect or estimate a false-positive rate. Losses under narrow/short positive inhibition prevent an assertion that the procedure reliably detects inhibition whenever it exists. Broad/long positive examples demonstrate success under those retained laws, with small test counts. The regimes have different laws and denominators; an aggregate success percentage would not estimate a meaningful common power. No inconvenient seed, cap or recovery history is replaced by a more favourable outcome.

The expanded matrix in Section~\ref{sec:s04-expanded} separates original LI, magnitude-scaled immigrant attenuation and total-intensity attenuation. Under the original unit-mark, immigrant-only generating law, magnitude-scaled immigrant-only inhibition has a positive mean gain in the broad/long regime, whereas total-intensity attenuation has a negative mean there. Both LI+ procedures are alternatives to that generating law, so this comparison does not measure their power under matched LI+ generation. The matrix reports Monte Carlo uncertainty by regime, and shows that gating can remove null overfitting while also selecting no inhibition in the narrow/short positive regime. This extends procedural sensitivity evidence without establishing physical absence or nominal test size.

The separate matched-generator extension in Section~\ref{sec:s04b-matched} generates magnitude-scaled immigrant-only and total-intensity histories. The generator counters change the interpretation of its narrow/short immigrant setting: 19 of 20 histories have no rejected proposals. This is a weakly realized control, not evidence of reliable recovery or its general impossibility. The total-attenuation setting rejects offspring and is a different experiment. Both matched raw procedures retain negative narrow/short mean gains against either floor. The original broad/long wins all use a fixed nuisance backbone. In the executed intercept sensitivity (Section~\ref{sec:s04b-refitted}), fitting the uninhibited background on training leaves matched mean gains of 5.1631 and 9.7367 nats, with 19/20 and 18/20 wins. This removes most of the large fixed-floor advantage but not all of it. The saved inhibitory parameters, gates and remaining nuisance coordinates are unchanged, so the comparison does not establish the performance of jointly refitted families. The fitted-parameter table also shows boundary concentration; an off-floor fit need not recover the generating parameters.

The completed magnitude-scaled observed comparison (Section~\ref{sec:s04-observed-completed}) retains the same conclusion: immigrant-only LI+ selects no inhibition in every fold. Raw total-intensity attenuation selects positive inhibition but loses held-out joint score in every fold; its complexity-gated version returns to the exact uninhibited floor. This fixed-backbone comparison addresses the revised diagnostic procedures, not a general test of physical stress shadows.

The completed conditional comparison (Section~\ref{sec:s05-completed}) retains discrepancies for both OC and EP under the two continuous-statistic omnibus procedures at each origin. The retained minimum-rank procedure does not reject at the declared level. Better comparative log score and conditional adequacy are therefore distinct: EP wins important score comparisons without thereby passing these diagnostic checks. Fitted states are held fixed, and no joint cross-origin null or parameter-uncertainty integration is claimed.

\section{Interpretation of the completed comparisons}

The New Zealand programme adaptively reused the same 612 test events: earlier test scores informed later model design and parameter bounds. Freezing each stage did not erase this dependence \citep{Dwork2015}. The inspected historical windows, overlapping training prefixes and dependent earthquake histories preclude event-iid confidence intervals and untouched-confirmation language. The singular zero-strength boundary does not come with a justified Wilks or half-chi-squared reference law. A bounded multistart search, independent score reconstruction and near-optimal forecast-spread check are useful controls, but are not global-optimization proofs.


The literature classifications in Appendix~\ref{app:literature_status} identify inherited triggering kernels, established stability machinery and restricted model-specific derivations. The contribution is the stated analysis and the executed procedure comparisons.

\section{Conclusion}

Mathematical legitimacy, identification, estimation, competitive prediction and calibration are distinct achievements. The finite-mode and finite-history results provide a restricted foundation; fitted nuisance parameters, stronger benchmarks, calibration and transfer constrain its empirical reach. The tested inhibition additions do not establish predictive success on the observed catalogues. This does not imply physical absence. Positive ablations, negative contrasts, boundary selections and completed failed gates remain part of that assessment.
